### Low-Rank Truncated-SVD Solver for Multi-Index ADP

Рассматривается глобальный шаг multi-index ADP при фиксированных локальных коэффициентах \(g_j\):

\[
\min_{B\in\mathbb R^{m\times d}}
F(B),
\]

\[
F(B)
=
\sum_{j=1}^{J}
N_j
\left\|
I_j-U_jB^\top g_j
\right\|_2^2
+
\lambda\|B-P\|_F^2,
\]

где

\[
P,B\in\mathbb R^{m\times d},\qquad
PP^\top=I_m,
\]

\[
g_j\in\mathbb R^m,\qquad
I_j\in\mathbb R^p,\qquad
U_j\in\mathbb R^{p\times d}.
\]

Требуется построить приближение ранга \(r<m\):

\[
\operatorname{rank}(B)\le r.
\]

Определим

\[
H_j=U_j^\top U_j,
\qquad
c_j=U_j^\top I_j.
\]

---

## 1. Rank-\(r\) параметризация

Ищем \(B\) в форме truncated SVD

\[
B_r=A\Sigma V^\top
=
\sum_{k=1}^{r}\sigma_k a_kv_k^\top,
\]

где

\[
A=[a_1,\ldots,a_r]\in\mathbb R^{m\times r},
\]

\[
V=[v_1,\ldots,v_r]\in\mathbb R^{d\times r},
\]

\[
A^\top A=I_r,
\qquad
V^\top V=I_r,
\]

\[
\Sigma=\operatorname{diag}(\sigma_1,\ldots,\sigma_r).
\]

Для центра \(j\)

\[
B_r^\top g_j
=
V\Sigma A^\top g_j.
\]

Обозначим

\[
\alpha_j=A^\top g_j.
\]

Тогда

\[
\begin{aligned}
\left\|
I_j-U_jV\Sigma\alpha_j
\right\|^2
={}&
\|I_j\|^2
-
2\alpha_j^\top\Sigma V^\top c_j
\\
&+
\alpha_j^\top
\Sigma V^\top H_jV\Sigma\alpha_j.
\end{aligned}
\]

Так как \(A\) и \(V\) имеют ортонормированные столбцы,

\[
\|A\Sigma V^\top\|_F^2
=
\|\Sigma\|_F^2
=
\sum_{k=1}^r\sigma_k^2.
\]

Кроме того,

\[
\|P\|_F^2=m.
\]

Следовательно,

\[
\boxed{
\|B_r-P\|_F^2
=
m+
\sum_{k=1}^r\sigma_k^2
-
2\operatorname{tr}
\left(
\Sigma A^\top PV
\right).
}
\]

---

## 2. Квадратичная задача по singular values

Для

\[
B_r=
\sum_{k=1}^r\sigma_ka_kv_k^\top
\]

имеем

\[
U_jB_r^\top g_j
=
\sum_{k=1}^{r}
\sigma_k(a_k^\top g_j)U_jv_k.
\]

Поэтому

\[
F(B_r)
=
C+
\sigma^\top M\sigma
-
2b^\top\sigma,
\]

где \(C\) не зависит от \(\sigma\), а

\[
\boxed{
M_{k\ell}
=
\sum_{j=1}^{J}
N_j
(a_k^\top g_j)
(a_\ell^\top g_j)
(U_jv_k)^\top(U_jv_\ell)
+
\lambda\delta_{k\ell}
}
\]

и

\[
\boxed{
b_k
=
\sum_{j=1}^{J}
N_j
(a_k^\top g_j)
I_j^\top U_jv_k
+
\lambda a_k^\top Pv_k.
}
\]

При фиксированных \(A,V\) оптимальные singular values находятся точно:

\[
\boxed{
M\sigma=b.
}
\]

То есть

\[
\sigma^\star=M^{-1}b
\]

при невырожденном \(M\). На практике следует решать систему через Cholesky/QR, а не вычислять \(M^{-1}\).

После решения отрицательный знак \(\sigma_k\) можно поглотить в \(a_k\) или \(v_k\).

---

## 3. Rank-1 подзадача

Для построения компонент последовательно рассмотрим текущее приближение \(B\) и добавку

\[
B^+=B+\sigma av^\top,
\]

где

\[
\|a\|_2=\|v\|_2=1.
\]

Определим текущие residuals

\[
e_j
=
I_j-U_jB^\top g_j.
\]

Тогда

\[
F(B+\sigma av^\top)-F(B)
=
D_B(a,v)\sigma^2
-
2R_B(a,v)\sigma,
\]

где

\[
\boxed{
D_B(a,v)
=
\sum_j
N_j
(a^\top g_j)^2
\|U_jv\|^2
+
\lambda
}
\]

и

\[
\boxed{
R_B(a,v)
=
\sum_j
N_j
(a^\top g_j)
v^\top U_j^\top e_j
+
\lambda a^\top(P-B)v.
}
\]

Для фиксированных \(a,v\)

\[
\boxed{
\sigma^\star(a,v)
=
\frac{R_B(a,v)}
     {D_B(a,v)}.
}
\]

Соответствующее уменьшение функционала равно

\[
\boxed{
F(B+\sigma^\star av^\top)-F(B)
=
-
\frac{R_B(a,v)^2}
     {D_B(a,v)}.
}
\]

Следовательно, оптимальная rank-1 компонента решает

\[
\boxed{
\max_{\|a\|=\|v\|=1}
\frac{R_B(a,v)^2}
     {D_B(a,v)}.
}
\]

---

## 4. Попеременное восстановление \(a\) и \(v\)

### Шаг по \(a\)

При фиксированном \(v\)

\[
R_B(a,v)=a^\top q_v,
\]

где

\[
q_v
=
\sum_j
N_j
g_j
\left(v^\top U_j^\top e_j\right)
+
\lambda(P-B)v.
\]

Кроме того,

\[
D_B(a,v)
=
a^\top G_v a,
\]

где

\[
\boxed{
G_v
=
\sum_j
N_j
\|U_jv\|^2
g_jg_j^\top
+
\lambda I_m.
}
\]

Поэтому

\[
\max_a
\frac{(a^\top q_v)^2}
     {a^\top G_va}
\]

имеет решение

\[
\boxed{
a
\propto
G_v^{-1}q_v.
}
\]

После решения выполняется нормировка

\[
a\leftarrow\frac{a}{\|a\|}.
\]

Матрица \(G_v\) имеет размер только \(m\times m\).

### Шаг по \(v\)

При фиксированном \(a\)

\[
R_B(a,v)=v^\top q_a,
\]

где

\[
q_a
=
\sum_j
N_j
(a^\top g_j)
U_j^\top e_j
+
\lambda(P-B)^\top a.
\]

Определим

\[
\boxed{
H_a
=
\sum_j
N_j
(a^\top g_j)^2
U_j^\top U_j
+
\lambda I_d.
}
\]

Тогда

\[
D_B(a,v)=v^\top H_av
\]

и

\[
\boxed{
v
\propto
H_a^{-1}q_a.
}
\]

После решения

\[
v\leftarrow\frac{v}{\|v\|}.
\]

Матрицу \(H_a\in\mathbb R^{d\times d}\) строить не требуется. Достаточно matrix-free операции

\[
H_ax
=
\sum_j
N_j
(a^\top g_j)^2
U_j^\top(U_jx)
+
\lambda x.
\]

Поэтому шаг по \(v\) можно решать CG/LSMR.

---

## 5. Последовательное rank-\(r\) восстановление

Инициализация:

\[
B_0=0.
\]

Для

\[
k=1,\ldots,r
\]

выполняется следующий цикл.

### Шаг 1. Residual

\[
e_j^{(k)}
=
I_j-U_jB_{k-1}^\top g_j.
\]

### Шаг 2. Инициализация нового \(v\)

Можно использовать случайный единичный вектор или ведущий правый singular vector residual-gradient

\[
Q_{k-1}
=
\sum_j
N_j
g_j
\left(U_j^\top e_j^{(k)}\right)^\top
+
\lambda(P-B_{k-1}).
\]

Полный \(Q_{k-1}\) строить не обязательно. Его действия на векторы вычисляются matrix-free.

### Шаг 3. Alternating rank-1 solve

До внутренней сходимости повторять

\[
q_v
=
\sum_j
N_jg_j
(v^\top U_j^\top e_j^{(k)})
+
\lambda(P-B_{k-1})v,
\]

\[
G_v
=
\sum_j
N_j
\|U_jv\|^2g_jg_j^\top
+
\lambda I_m,
\]

\[
a
\leftarrow
\operatorname{normalize}
\left(G_v^{-1}q_v\right),
\]

после чего

\[
q_a
=
\sum_j
N_j
(a^\top g_j)
U_j^\top e_j^{(k)}
+
\lambda(P-B_{k-1})^\top a,
\]

\[
H_a x
=
\sum_j
N_j
(a^\top g_j)^2
U_j^\top(U_jx)
+
\lambda x,
\]

\[
v
\leftarrow
\operatorname{normalize}
\left(H_a^{-1}q_a\right).
\]

### Шаг 4. Оптимальный масштаб

Вычислить

\[
R=
R_{B_{k-1}}(a,v),
\qquad
D=
D_{B_{k-1}}(a,v),
\]

и

\[
\boxed{
\sigma_k=\frac RD.
}
\]

### Шаг 5. Rank-1 update

\[
\widetilde B_k
=
B_{k-1}
+
\sigma_kav^\top.
\]

### Шаг 6. Compact SVD

Выполнить

\[
\widetilde B_k
=
A_k\Sigma_kV_k^\top
\]

и оставить не более \(k\) ненулевых компонент.

Полный SVD матрицы \(m\times d\) не требуется.

Если

\[
\widetilde B_k=LR^\top,
\qquad
L\in\mathbb R^{m\times k},
\quad
R\in\mathbb R^{d\times k},
\]

выполняются thin QR

\[
L=Q_LR_L,
\qquad
R=Q_RR_R
\]

и SVD только матрицы

\[
R_LR_R^\top\in\mathbb R^{k\times k}.
\]

---

## 6. Совместная переоценка singular values

После compact SVD желательно не сохранять жадно найденные \(\sigma_1,\ldots,\sigma_k\), а решить точную условную задачу.

Для текущих

\[
A_k=[a_1,\ldots,a_k],
\qquad
V_k=[v_1,\ldots,v_k]
\]

строятся

\[
(M_k)_{s\ell}
=
\sum_j
N_j
(a_s^\top g_j)
(a_\ell^\top g_j)
(U_jv_s)^\top(U_jv_\ell)
+
\lambda\delta_{s\ell},
\]

\[
(b_k)_s
=
\sum_j
N_j
(a_s^\top g_j)
I_j^\top U_jv_s
+
\lambda a_s^\top Pv_s.
\]

После чего

\[
\boxed{
M_k\sigma=b_k.
}
\]

Затем

\[
B_k
=
A_k\operatorname{diag}(\sigma)V_k^\top.
\]

Этот шаг оптимизирует все \(k\) singular values одновременно и устраняет часть ошибки жадного восстановления.

---

## 7. Критерий остановки

Алгоритм можно остановить до достижения заданного \(r\), если новая компонента практически не уменьшает функционал:

\[
\frac{R^2/D}
     {\max(1,F(B_{k-1}))}
<
\varepsilon_{\mathrm{rank}}.
\]

Дополнительный критерий определяется singular spectrum:

\[
\frac{\sigma_k^2}
{\sum_{\ell=1}^{k}\sigma_\ell^2}
<
\varepsilon_{\mathrm{sv}}.
\]

Поэтому effective rank может определяться адаптивно.

---

## 8. Полный алгоритм

**Вход:**  
\(\{I_j,U_j,N_j,g_j\}_{j=1}^J\), \(P\), \(\lambda\), максимальный ранг \(r\), допуски \(\varepsilon_{\mathrm{inner}},\varepsilon_{\mathrm{rank}}\).

**Инициализация:**

\[
B_0=0.
\]

**Для \(k=1,\ldots,r\):**

1. Вычислить
   \[
   e_j=I_j-U_jB_{k-1}^\top g_j.
   \]

2. Инициализировать \(v\).

3. Повторять до внутренней сходимости:
   \[
   a\leftarrow
   \operatorname{normalize}(G_v^{-1}q_v),
   \]
   \[
   v\leftarrow
   \operatorname{normalize}(H_a^{-1}q_a).
   \]

4. Вычислить
   \[
   \sigma=R_B(a,v)/D_B(a,v).
   \]

5. Если
   \[
   R_B(a,v)^2/D_B(a,v)
   \]
   меньше порога, остановиться.

6. Выполнить
   \[
   B_k\leftarrow B_{k-1}+\sigma av^\top.
   \]

7. Сделать compact SVD \(B_k\).

8. Для полученных \(A_k,V_k\) решить
   \[
   M_k\sigma=b_k.
   \]

9. Обновить
   \[
   B_k=A_k\operatorname{diag}(\sigma)V_k^\top.
   \]

**Выход:**

\[
\boxed{
B_r=A_r\Sigma_rV_r^\top.
}
\]

---

## 9. Интерпретация как generalized truncated SVD

Определим оператор

\[
\boxed{
\mathcal L(X)
=
\sum_j
N_j
(g_jg_j^\top)
X
(U_j^\top U_j)
+
\lambda X
}
\]

и

\[
\boxed{
C_B
=
\sum_j
N_jg_jc_j^\top+\lambda P.
}
\]

Полное unconstrained решение \(B_\star\) удовлетворяет

\[
\mathcal L(B_\star)=C_B.
\]

Поскольку функционал квадратичен,

\[
F(B)
=
F(B_\star)
+
\left\langle
B-B_\star,
\mathcal L(B-B_\star)
\right\rangle.
\]

Следовательно, точная rank-\(r\) задача имеет форму

\[
\boxed{
B_r^\star
=
\arg\min_{\operatorname{rank}(B)\le r}
\|B-B_\star\|_{\mathcal L}^2
}
\]

с

\[
\|X\|_{\mathcal L}^2
=
\langle X,\mathcal L(X)\rangle.
\]

Поэтому предложенный алгоритм можно рассматривать как последовательное вычисление **generalized singular components в Hessian-метрике HPAO**.

Обычный truncated SVD является частным случаем, когда

\[
\mathcal L(X)\propto X.
\]

---

## 10. Изотропное приближение

Если

\[
U_j^\top U_j
\approx
\eta_jI_d,
\]

то

\[
\mathcal L(B)
\approx
GB,
\]

где

\[
G
=
\sum_j
N_j\eta_jg_jg_j^\top
+
\lambda I_m.
\]

Тогда

\[
F(B)
=
\operatorname{const}
+
\left\|
G^{1/2}B-G^{-1/2}C_B
\right\|_F^2.
\]

С заменой

\[
Z=G^{1/2}B
\]

rank сохраняется, и по теореме Eckart–Young

\[
Z_r
=
\left[
G^{-1/2}C_B
\right]_r.
\]

Следовательно,

\[
\boxed{
B_r
=
G^{-1/2}
\left[
G^{-1/2}C_B
\right]_r.
}
\]

Эта формула задаёт прямой approximate solver:

\[
\boxed{
G,C_B
\rightarrow
G^{-1/2}C_B
\rightarrow
\text{truncated SVD}_r
\rightarrow
B_r.
}
\]

Её можно использовать либо как самостоятельную аппроксимацию, либо как инициализацию общего low-rank solver.

**Замечание.** Ограничение \(\operatorname{rank}(B)\le r<m\) означает, что этот solver намеренно аппроксимирует полный HPAO-шаг более низкоранговой матрицей. Поэтому после получения \(B_r\) EDR-базис следует восстанавливать тем же спектральным шагом, что и в полном алгоритме; \(B_r\) не следует непосредственно интерпретировать как новый \(m\)-мерный ортонормированный базис.
